% ============================================================
%  Paper 02: Topology of the 24-Cell and F4 Symmetries in Quantum Field Theory
%  Target: RevTeX 4-2 Compilation (SDGFT Series)
% ============================================================
\documentclass[amsmath,amssymb,aps,prd,twocolumn,superscriptaddress,nofootinbib]{revtex4-2}

\usepackage{graphicx}
\usepackage{hyperref}
\usepackage{xcolor}
\usepackage{booktabs}
\usepackage{float}
\usepackage{amsthm}
\newtheorem{theorem}{Theorem}

% ── SDGFT shared macros ──────────────────────────────────────
\usepackage{sdgft-macros}

% ── Document ─────────────────────────────────────────────────
\begin{document}

\preprint{SDGFT-2026-02}

\title{Topology of the 24-Cell and F4 Symmetries in Quantum Field Theory}

\author{David A. Besemer}
\email{david.besemer@sdgft.org}
\affiliation{Independent Researcher}

\date{\today}

\begin{abstract}
We investigate the algebraic structure of the 24-cell polytope {3,4,3} and its automorphism group F4. By decomposing the 24-cell into three orthogonal 16-cell sub-lattices, we demonstrate the emergence of D4 triality. We establish how the gauge symmetries of the Standard Model arise naturally from Dynkin diagram deletions of the F4 root system, presenting a top-down geometric derivation of gauge groups in quantum field theory.
\end{abstract}

\maketitle

\section{Introduction}
The unification of gauge symmetries with gravity is a long-standing goal of mathematical physics. We approach this by studying the Dynkin diagram of F4, the exceptional Lie algebra associated with the symmetries of the 24-cell.

\section{Theoretical Formulation}
Here we formulate the core mathematical relations governing the system:
\begin{equation}
  24 = 8 + 8 + 8 \quad (\text{decomposition into three orthogonal 16-cells})
\end{equation}
\begin{equation}
  F_4 \xrightarrow{\text{deletion}} SU(3)_C \times SU(2)_L \times U(1)_Y
\end{equation}
\section{Dynkin Deletions}
By removing peripheral Dynkin nodes of the F4 algebra, we isolate the subalgebra corresponding to the Standard Model gauge group, demonstrating that flavour triality is topological.

\section{Conclusion}
The root system of the 24-cell uniquely determines the gauge structure of the early universe, precluding any additional dark gauge forces at the Planck scale.



\begin{thebibliography}{99}
\bibitem{Goldstein_Paper01} D. A. Besemer, ``Axiomatic Foundations of SDGFT,'' SDGFT-2026-01 (in preparation).
\bibitem{Goldstein_Paper04} D. A. Besemer, ``Effective Spectral Dimension and the Banach Fixed-Point Theorem in SDGFT,'' SDGFT-2026-04.
\bibitem{Goldstein_Code} D. A. Besemer, \texttt{sdgft} Python package, \url{https://github.com/david-besemer/sdgft} (2026).
\end{thebibliography}

\end{document}
